% ============================================================
%  CAHIER DES CHARGES — APPLICATION CAYA
%  Club des Amis de Yaoundé — Dématérialisation Complète
%  Version 1.0 — Mars 2026
% ============================================================
\documentclass[12pt, a4paper]{report}

% ---------- Encodage & Langue ----------
\usepackage[utf8]{inputenc}
\usepackage[T1]{fontenc}
\usepackage[french]{babel}

% ---------- Mise en page ----------
\usepackage[top=2.5cm, bottom=2.5cm, left=2.8cm, right=2.8cm]{geometry}
\usepackage{setspace}
\onehalfspacing

% ---------- Polices ----------
\usepackage{lmodern}
\usepackage{microtype}

% ---------- Couleurs & Graphiques ----------
\usepackage[dvipsnames, table]{xcolor}
\usepackage{graphicx}
\usepackage{tikz}

% Palette CAYA
\definecolor{cayaBlue}{RGB}{0, 71, 130}
\definecolor{cayaGold}{RGB}{201, 155, 0}
\definecolor{cayaLight}{RGB}{230, 241, 255}
\definecolor{cayaGray}{RGB}{245, 245, 245}
\definecolor{warningRed}{RGB}{200, 40, 40}
\definecolor{successGreen}{RGB}{30, 130, 60}

% ---------- En-têtes & pieds de page ----------
\usepackage{fancyhdr}
\pagestyle{fancy}
\fancyhf{}
\fancyhead[L]{\small\textcolor{cayaBlue}{\textbf{CAYA} — Cahier des Charges}}
\fancyhead[R]{\small\textcolor{cayaGray!50!black}{v1.0 · Mars 2026}}
\fancyfoot[C]{\small\thepage}
\renewcommand{\headrulewidth}{0.4pt}
\renewcommand{\footrulewidth}{0pt}
% ---------- Tableaux ----------
\usepackage{booktabs}
\usepackage{longtable}
\usepackage{tabularx}
\usepackage{multirow}
\usepackage{array}
\usepackage{colortbl}
\newcolumntype{L}[1]{>{\raggedright\arraybackslash}p{#1}}
\newcolumntype{C}[1]{>{\centering\arraybackslash}p{#1}}
\newcolumntype{R}[1]{>{\raggedleft\arraybackslash}p{#1}}

% ---------- Listes ----------
\usepackage{enumitem}
\setlist[itemize]{leftmargin=1.5em, itemsep=2pt}
\setlist[enumerate]{leftmargin=1.5em, itemsep=2pt}

% ---------- Boîtes ----------
\usepackage{tcolorbox}
\tcbuselibrary{breakable, skins, listings}

\newtcolorbox{ruleBox}[1]{
  colback=cayaLight, colframe=cayaBlue,
  fonttitle=\bfseries\small, title=#1,
  breakable, boxrule=0.8pt, arc=3pt
}

\newtcolorbox{invariantBox}[1]{
  colback=cayaGold!10, colframe=cayaGold!80!black,
  fonttitle=\bfseries\small\color{cayaGold!80!black},
  title=Invariant : #1,
  breakable, boxrule=0.8pt, arc=3pt
}

\newtcolorbox{warningBox}{
  colback=warningRed!8, colframe=warningRed,
  fonttitle=\bfseries\small\color{warningRed},
  title=Attention,
  breakable, boxrule=0.8pt, arc=3pt
}

\newtcolorbox{codeBox}{
  colback=gray!8, colframe=gray!50,
  fontupper=\small\ttfamily,
  breakable, boxrule=0.5pt, arc=2pt
}

% ---------- Code ----------
\usepackage{listings}
\lstset{
  basicstyle=\ttfamily\footnotesize,
  backgroundcolor=\color{gray!8},
  frame=single, framerule=0.4pt,
  breaklines=true,
  keepspaces=true,
  showstringspaces=false
}

% ---------- Hyper-liens ----------
\usepackage[hidelinks, colorlinks=true,
  linkcolor=cayaBlue, urlcolor=cayaBlue, citecolor=cayaBlue]{hyperref}
\usepackage{bookmark}

% ---------- Utilitaires ----------
\usepackage{ifthen}
\usepackage{calc}
\usepackage{xspace}
\usepackage{soul}         % surlignage
\usepackage{pifont}       % symboles
\newcommand{\cmark}{\textcolor{successGreen}{\ding{51}}}
\newcommand{\xmark}{\textcolor{warningRed}{\ding{55}}}

% ---------- Titre du document ----------
\title{%
  \vspace{-1cm}
  {\Huge\bfseries\color{cayaBlue} Cahier des Charges}\\[0.4cm]
  {\large\bfseries Application de Gestion\\[0.1cm]
  \textbf{\color{cayaGold}CLUB DES AMIS DE YAOUNDÉ (CAYA)}}\\[0.3cm]
  {\normalsize Dématérialisation Complète du Système de Tontine}\\[0.5cm]
  \rule{0.6\textwidth}{1pt}\\[0.4cm]
  {\small Version 1.0 \quad|\quad Mars 2026 \quad|\quad Confidentiel}
}
\author{}
\date{}

% ============================================================
\begin{document}
% ============================================================

\maketitle
\thispagestyle{empty}

\vfill
\begin{center}
  \begin{tabular}{ll}
    \toprule
    \textbf{Projet}    & Application CAYA \\
    \textbf{Version}   & 1.0 \\
    \textbf{Date}      & Mars 2026 \\
    \textbf{Statut}    & En modélisation \\
    \textbf{Équipe}    & Ingénierie Logicielle \\
    \bottomrule
  \end{tabular}
\end{center}

\newpage
\tableofcontents
\newpage

% ============================================================
\chapter*{Résumé Exécutif}
\addcontentsline{toc}{chapter}{Résumé Exécutif}
% ============================================================

Le présent document constitue le cahier des charges complet de l'application de gestion du
\textbf{Club des Amis de Yaoundé (CAYA)}, une association d'épargne rotative fondée en 2000 à Yaoundé.

\medskip
Ce cahier formalise la \textbf{dématérialisation intégrale} d'un système de gestion tontine jusqu'ici
géré manuellement sous tableur Excel. L'objectif est de produire une application robuste, modulaire,
conçue pour fonctionner sur de multiples exercices fiscaux consécutifs sans perte d'intégrité des données.

\medskip
\begin{center}
\rowcolors{2}{cayaLight}{white}
\begin{tabular}{L{5cm}L{8cm}}
  \toprule
  \rowcolor{cayaBlue}
  \textcolor{white}{\textbf{Paramètre}} & \textcolor{white}{\textbf{Valeur}} \\
  \midrule
  Nom officiel      & Club des Amis de Yaoundé \\
  Acronyme          & CAYA \\
  Fondé en          & 2000 \\
  Siège             & Yaoundé, Cameroun \\
  Devise            & Paix – Solidarité – Épanouissement \\
  Membres actifs    & 32 (exercice 2025-2026) \\
  Capital courant   & 10 000 000 XAF \\
  Valeur d'une part & 100 000 XAF \\
  Taux prêts        & 4\% / mois (encours glissant) \\
  Exercice type     & 12 mois (Août $\rightarrow$ Juillet) \\
  \bottomrule
\end{tabular}
\end{center}

% ============================================================
\chapter{Contexte et Objectifs}
% ============================================================

\section{Contexte}

Le CAYA est une tontine améliorée regroupant une trentaine de membres. Chaque exercice de 12 mois,
les membres effectuent des dépôts mensuels (parts, épargne, secours, pot), peuvent contracter des
prêts auprès de la caisse commune, et perçoivent à la \textit{cassation} leur épargne augmentée
des intérêts générés par les prêts redistribués proportionnellement.

\medskip
Jusqu'à ce jour, l'ensemble des calculs financiers — intérêts sur encours, redistribution proportionnelle,
suivi des remboursements, bilan de cassation — est réalisé manuellement dans un fichier Excel
multi-feuilles. Cette approche génère des \textbf{incohérences documentées} (écarts de 532 400 XAF
entre deux versions de calcul, erreurs de recopie manuelle créant des dettes fantômes).

\section{Objectifs du Projet}

\begin{enumerate}[label=\textbf{O\arabic*.}]
  \item \textbf{Automatisation intégrale} des calculs financiers (intérêts, redistribution, cassation).
  \item \textbf{Élimination des saisies manuelles redondantes} — toute donnée calculable doit être dérivée.
  \item \textbf{Traçabilité complète} — chaque opération financière est auditée, horodatée, attribuée.
  \item \textbf{Pérennité multi-exercices} — le système doit fonctionner de manière continue sur des années.
  \item \textbf{Gouvernance renforcée} — le contrôle d'accès par rôle garantit que chaque acteur n'effectue
        que les opérations qui lui sont autorisées.
  \item \textbf{Exportabilité} — les données doivent être exportables en PDF et Excel pour les réunions.
\end{enumerate}

\section{Périmètre}

L'application couvre \textbf{une seule tontine} : le CAYA. Elle n'est pas conçue pour gérer plusieurs
organisations simultanément. La conception est modulaire afin de garantir l'absence de régressions
lors de l'ajout de fonctionnalités.

% ============================================================
\chapter{Architecture Générale}
% ============================================================

\section{Principes Architecturaux}

L'application repose sur les principes suivants :

\begin{itemize}
  \item \textbf{Domain-Driven Design (DDD)} — les agrégats correspondent aux invariants métier.
  \item \textbf{Architecture en couches} — Présentation / Application / Domaine / Infrastructure.
  \item \textbf{CQRS léger} — séparation des opérations de lecture (vues rapides) et d'écriture (commandes métier).
  \item \textbf{Event sourcing partiel} — l'audit log et les ledgers constituent une trace événementielle immuable.
  \item \textbf{Conception module par module} — chaque module est développé, testé et validé isolément.
\end{itemize}

\section{Stack Technique}

\begin{center}
\rowcolors{2}{cayaLight}{white}
\begin{tabular}{L{3.5cm}L{4.5cm}L{5cm}}
  \toprule
  \rowcolor{cayaBlue}
  \textcolor{white}{\textbf{Couche}} &
  \textcolor{white}{\textbf{Technologie}} &
  \textcolor{white}{\textbf{Rôle}} \\
  \midrule
  Backend API      & NestJS (TypeScript)   & API REST sécurisée, logique métier \\
  Base de données  & MySQL 8.0             & Persistance, triggers, procédures stockées \\
  ORM              & Prisma                & Schéma typé, migrations \\
  Frontend Web     & Next.js / React       & Interface bureau (Trésorier, Bureau) \\
  Mobile           & Flutter (Dart)        & Interface membres — iOS \& Android \\
  Authentification & JWT + Refresh Token   & Auth phone/username, sessions longues \\
  Notifications    & WhatsApp API + SMS    & Rappels, confirmations \\
  \bottomrule
\end{tabular}
\end{center}

\section{Décomposition en Modules}

\begin{center}
\rowcolors{2}{cayaLight}{white}
\begin{tabular}{C{1cm}L{4cm}L{8cm}}
  \toprule
  \rowcolor{cayaBlue}
  \textcolor{white}{\textbf{\#}} &
  \textcolor{white}{\textbf{Module}} &
  \textcolor{white}{\textbf{Responsabilité}} \\
  \midrule
  01 & Auth \& RBAC           & Authentification, rôles, sessions, audit log \\
  02 & Configuration          & Paramètres CAYA, snapshot exercice, événements secours \\
  03 & Membres                & Profils, contacts d'urgence, parrainage \\
  04 & Exercice Fiscal        & Cycle de vie exercice, adhésions, engagements de parts \\
  05 & Sessions \& Transactions & Réunions mensuelles, pointage, référencement \\
  06 & Épargne \& Intérêts    & Ledgers épargne, redistribution mensuelle \\
  07 & Prêts                  & Octroi, accrual mensuel, remboursements, report \\
  08 & Caisse de Secours      & Fonds, événements, positions membres \\
  09 & Bénéficiaires          & Planning de la « bouffe », désignation, livraison \\
  10 & Cassation \& Rollover  & Clôture d'exercice, redistribution finale, report prêts \\
  11 & Rapports \& Exports    & PDF, Excel, graphiques, tableau de bord \\
  12 & Notifications          & WhatsApp, SMS, rappels automatiques \\
  \bottomrule
\end{tabular}
\end{center}

% ============================================================
\chapter{Modèle de Données — Entités et Règles Métier}
% ============================================================

\section{Module 01 — Authentification et Contrôle d'Accès}

\subsection{Entité \texttt{User}}

\begin{codeBox}
User {
  id            : ShortUUID(10)   PK, auto-généré
  username      : VARCHAR(50)     UNIQUE, NOT NULL
  phone         : VARCHAR(20)     UNIQUE, NOT NULL
  passwordHash  : VARCHAR(255)    NOT NULL
  role          : BureauRole      NOT NULL, DEFAULT MEMBRE
  isActive      : BOOLEAN         DEFAULT true
  createdAt     : TIMESTAMP       DEFAULT now()
  updatedAt     : TIMESTAMP       auto-updated
  lastLoginAt   : TIMESTAMP       nullable
}
\end{codeBox}

\subsection{Entité \texttt{RefreshToken}}

\begin{codeBox}
RefreshToken {
  id          : ShortUUID(10)   PK
  userId      : ShortUUID       FK → User
  tokenHash   : VARCHAR(255)    UNIQUE
  expiresAt   : TIMESTAMP
  revokedAt   : TIMESTAMP       nullable
  createdAt   : TIMESTAMP
  userAgent   : VARCHAR(500)    nullable
}
\end{codeBox}

\subsection{Énumération \texttt{BureauRole}}

\begin{center}
\rowcolors{2}{cayaLight}{white}
\begin{tabular}{L{4.5cm}L{9cm}}
  \toprule
  \rowcolor{cayaBlue}
  \textcolor{white}{\textbf{Rôle}} &
  \textcolor{white}{\textbf{Responsabilités principales}} \\
  \midrule
  \texttt{SUPER\_ADMIN}             & Modification config en cours d'exercice, accès total \\
  \texttt{PRESIDENT}                & Autorisation secours, désignation bénéficiaires, validation sessions \\
  \texttt{VICE\_PRESIDENT}          & Délégation Président, autorisation secours \\
  \texttt{TRESORIER}                & Saisie des transactions, déblocage prêts, exécution cassation \\
  \texttt{TRESORIER\_ADJOINT}       & Assistance Trésorier \\
  \texttt{SECRETAIRE\_GENERAL}      & Accès audit log, gestion réunions \\
  \texttt{SECRETAIRE\_ADJOINT}      & Assistance Secrétaire \\
  \texttt{COMMISSAIRE\_AUX\_COMPTES}& Lecture financière complète, pas de saisie \\
  \texttt{MEMBRE}                   & Consultation de sa propre fiche uniquement \\
  \bottomrule
\end{tabular}
\end{center}

\begin{invariantBox}{AUTH-01 à AUTH-04}
\begin{itemize}
  \item \textbf{AUTH-01} : Un seul \texttt{User} avec \texttt{role = PRESIDENT} et \texttt{isActive = true}.
  \item \textbf{AUTH-02} : Un seul \texttt{User} avec \texttt{role = VICE\_PRESIDENT} et \texttt{isActive = true}.
  \item \textbf{AUTH-03} : \texttt{phone} de \texttt{User} est synchronisé avec \texttt{MemberProfile.phone1}.
  \item \textbf{AUTH-04} : Tout changement de rôle génère une entrée dans \texttt{AuditLog}.
\end{itemize}
\end{invariantBox}

\subsection{Matrice RBAC}

\begin{center}
\small
\rowcolors{2}{cayaLight}{white}
\begin{tabular}{L{3.8cm}C{1.5cm}C{1.5cm}C{1.8cm}C{1.8cm}C{1.8cm}}
  \toprule
  \rowcolor{cayaBlue}
  \textcolor{white}{\textbf{Permission}} &
  \textcolor{white}{\textbf{S.Admin}} &
  \textcolor{white}{\textbf{Prés./VP}} &
  \textcolor{white}{\textbf{Trés.}} &
  \textcolor{white}{\textbf{Secr.}} &
  \textcolor{white}{\textbf{Membre}} \\
  \midrule
  Modifier config exercice actif   & \cmark & \xmark & \xmark & \xmark & \xmark \\
  Lire configuration               & \cmark & \cmark & \cmark & \cmark & \xmark \\
  Saisir transactions session      & \cmark & \cmark & \cmark & \xmark & \xmark \\
  Autoriser prêt                   & \cmark & \cmark & \xmark & \xmark & \xmark \\
  Débloquer prêt (décaissement)    & \cmark & \xmark & \cmark & \xmark & \xmark \\
  Autoriser secours                & \cmark & \cmark & \xmark & \xmark & \xmark \\
  Désigner bénéficiaire            & \cmark & \cmark & \xmark & \xmark & \xmark \\
  Exécuter cassation               & \cmark & \xmark & \cmark & \xmark & \xmark \\
  Rouvrir session clôturée         & \cmark & \xmark & \xmark & \xmark & \xmark \\
  Consulter audit log              & \cmark & \cmark & \xmark & \cmark & \xmark \\
  Consulter sa fiche               & \cmark & \cmark & \cmark & \cmark & \cmark \\
  \bottomrule
\end{tabular}
\end{center}

% -------------------------------------------------------
\section{Module 02 — Configuration}

\subsection{Entité \texttt{TontineConfig} (Singleton)}

\begin{codeBox}
TontineConfig {
  id                      : VARCHAR    PK DEFAULT 'caya'
  -- Identité --
  name                    : VARCHAR(255)
  acronym                 : VARCHAR(20)
  foundedYear             : INT
  motto                   : VARCHAR(500)
  headquartersCity        : VARCHAR(255)
  registrationNumber      : VARCHAR(100)  nullable
  -- Parts --
  shareUnitAmount         : DECIMAL(15,2)  DEFAULT 100000
  halfShareAmount         : DECIMAL(15,2)  DEFAULT 50000
  potMonthlyAmount        : DECIMAL(15,2)  DEFAULT 3000
  maxSharesPerMember      : INT            DEFAULT 10
  -- Épargne --
  mandatoryInitialSavings : DECIMAL(15,2)  DEFAULT 100000
  -- Prêts --
  loanMonthlyRate         : DECIMAL(5,4)   DEFAULT 0.04
  minLoanAmount           : DECIMAL(15,2)  DEFAULT 10000
  maxLoanAmount           : DECIMAL(15,2)  DEFAULT 1000000
  maxLoanMultiplier       : INT            DEFAULT 5
  minSavingsToLoan        : DECIMAL(15,2)  DEFAULT 100000
  maxConcurrentLoans      : INT            DEFAULT 2
  -- Secours --
  rescueFundTarget        : DECIMAL(15,2)  DEFAULT 50000
  rescueFundMinBalance    : DECIMAL(15,2)  DEFAULT 25000
  -- Inscription --
  registrationFeeNew      : DECIMAL(15,2)
  registrationFeeReturning: DECIMAL(15,2)
  -- Méta --
  updatedAt               : TIMESTAMP
  updatedById             : ShortUUID FK → User
}
\end{codeBox}

\subsection{Entité \texttt{FiscalYearConfig} (Snapshot Immuable)}

\begin{codeBox}
FiscalYearConfig {
  id                  : ShortUUID(10)  PK
  fiscalYearId        : ShortUUID      FK → FiscalYear  UNIQUE
  snapshotAt          : TIMESTAMP
  snapshotById        : ShortUUID      FK → User
  -- Copie gelée des règles financières (identique à TontineConfig) --
  shareUnitAmount     : DECIMAL(15,2)
  loanMonthlyRate     : DECIMAL(5,4)
  maxLoanMultiplier   : INT
  minSavingsToLoan    : DECIMAL(15,2)
  maxConcurrentLoans  : INT
  rescueFundTarget    : DECIMAL(15,2)
  rescueFundMinBalance: DECIMAL(15,2)
  registrationFeeNew  : DECIMAL(15,2)
  registrationFeeReturning : DECIMAL(15,2)
  interestPoolMethod  : InterestPoolMethod
  -- Exception SUPER_ADMIN uniquement --
  forcedModifiedAt    : TIMESTAMP    nullable
  forcedModifiedById  : ShortUUID    nullable
  forcedModifiedReason: TEXT         nullable
}

InterestPoolMethod (ENUM) :
  THEORETICAL  -- encours[m-1] × 0.04, indépendant des paiements
  ACTUAL       -- ∑ RBT_INTEREST effectivement encaissés ce mois
\end{codeBox}

\begin{invariantBox}{CONF-01 à CONF-03}
\begin{itemize}
  \item \textbf{CONF-01} : \texttt{FiscalYearConfig} est créé lors de l'activation de l'exercice (\texttt{PENDING → ACTIVE})
        par snapshot de \texttt{TontineConfig} — immuable ensuite.
  \item \textbf{CONF-02} : Toute modification post-activation requiert \texttt{role = SUPER\_ADMIN}
        et génère une entrée \texttt{AuditLog} avec \texttt{forcedModifiedReason} obligatoire.
  \item \textbf{CONF-03} : \texttt{interestPoolMethod} est switchable par exercice, mais uniquement
        en statut \texttt{PENDING}.
\end{itemize}
\end{invariantBox}

\subsection{Montants de Secours par Événement}

\begin{center}
\rowcolors{2}{cayaLight}{white}
\begin{tabular}{L{6cm}R{4cm}}
  \toprule
  \rowcolor{cayaBlue}
  \textcolor{white}{\textbf{Événement}} &
  \textcolor{white}{\textbf{Montant (XAF)}} \\
  \midrule
  Décès d'un membre              & 300 000 \\
  Décès d'un parent proche       & 100 000 \\
  Mariage                        & 200 000 \\
  Naissance                      & 25 000 \\
  Maladie grave                  & 50 000 \\
  Promotion professionnelle      & 50 000 \\
  \bottomrule
\end{tabular}
\end{center}

% -------------------------------------------------------
\section{Module 03 — Membres}

\subsection{Entité \texttt{MemberProfile}}

\begin{codeBox}
MemberProfile {
  id                : ShortUUID(10)  PK
  memberCode        : VARCHAR(10)    UNIQUE  -- format CAYA-00042
  userId            : ShortUUID      FK → User  UNIQUE
  firstName         : VARCHAR(100)   NOT NULL
  lastName          : VARCHAR(100)   NOT NULL
  phone1            : VARCHAR(20)    NOT NULL
  phone2            : VARCHAR(20)    nullable
  attachmentPhoto   : BYTEA          nullable  -- JPEG compressé < 100 KB
  neighborhood      : VARCHAR(255)   NOT NULL
  locationDetail    : TEXT           nullable
  mobileMoneyType   : VARCHAR(50)    nullable  -- MTN / Orange
  mobileMoneyNumber : VARCHAR(20)    nullable
  sponsorId         : ShortUUID      FK → MemberProfile  nullable
  createdAt         : TIMESTAMP      DEFAULT now()
  updatedAt         : TIMESTAMP      auto-updated
}
\end{codeBox}

\subsection{Entité \texttt{EmergencyContact}}

\begin{codeBox}
EmergencyContact {
  id          : ShortUUID(10)  PK
  profileId   : ShortUUID      FK → MemberProfile
  fullName    : VARCHAR(255)   NOT NULL
  phone       : VARCHAR(20)    NOT NULL
  relation    : VARCHAR(100)   nullable
}
\end{codeBox}

% -------------------------------------------------------
\section{Module 04 — Exercice Fiscal et Adhésions}

\subsection{Entité \texttt{FiscalYear}}

\begin{codeBox}
FiscalYear {
  id              : ShortUUID(10)       PK
  label           : VARCHAR(20)         UNIQUE  -- ex: "2025-2026"
  startDate       : DATE                NOT NULL
  endDate         : DATE                NOT NULL
  cassationDate   : DATE                NOT NULL
  loanDueDate     : DATE                NOT NULL  -- avant cassationDate
  status          : FiscalYearStatus    DEFAULT PENDING
  openedAt        : TIMESTAMP           nullable
  openedById      : ShortUUID FK → User nullable
  closedAt        : TIMESTAMP           nullable
  closedById      : ShortUUID FK → User nullable
  notes           : TEXT                nullable
}

FiscalYearStatus (ENUM) : PENDING → ACTIVE → CASSATION → CLOSED → ARCHIVED

CHECK : startDate < loanDueDate < cassationDate ≤ endDate
TRIGGER : Aucun chevauchement de dates avec un FiscalYear existant
\end{codeBox}

\subsection{Machine d'États — \texttt{FiscalYear}}

\begin{center}
\begin{tikzpicture}[
  node distance=2.8cm,
  state/.style={rectangle, rounded corners=5pt, draw=cayaBlue,
                fill=cayaLight, minimum width=2.8cm, minimum height=0.9cm,
                font=\small\bfseries},
  arrow/.style={->, thick, cayaBlue}
]
  \node[state] (P) {PENDING};
  \node[state, right of=P, xshift=0.5cm] (A) {ACTIVE};
  \node[state, right of=A, xshift=0.5cm] (C) {CASSATION};
  \node[state, right of=C, xshift=0.5cm] (CL) {CLOSED};
  \node[state, below of=CL] (AR) {ARCHIVED};

  \draw[arrow] (P) -- node[above, font=\tiny]{activate()} (A);
  \draw[arrow] (A) -- node[above, font=\tiny]{openCassation()} (C);
  \draw[arrow] (C) -- node[above, font=\tiny]{close()} (CL);
  \draw[arrow] (CL) -- node[right, font=\tiny]{archive()} (AR);
\end{tikzpicture}
\end{center}

\subsection{Entité \texttt{Membership}}

\begin{codeBox}
Membership {
  id                  : ShortUUID(10)   PK
  profileId           : ShortUUID       FK → MemberProfile
  fiscalYearId        : ShortUUID       FK → FiscalYear
  status              : MemberStatus    DEFAULT ACTIVE
  joinedAt            : DATE            NOT NULL
  joinedAtMonth       : INT             NOT NULL  -- mois 1..12
  enrollmentType      : EnrollmentType
  catchUpAmount       : DECIMAL(15,2)   DEFAULT 0
  catchUpPaid         : BOOLEAN         DEFAULT false
  registrationFeePaid : BOOLEAN         DEFAULT false
  rescueContribPaid   : BOOLEAN         DEFAULT false
  initialSavingsPaid  : BOOLEAN         DEFAULT false
  UNIQUE(profileId, fiscalYearId)
}

MemberStatus   : ACTIVE | INACTIVE | SUSPENDED | EXCLUDED
EnrollmentType : NEW | RETURNING | MID_YEAR
\end{codeBox}

\begin{ruleBox}{Règle d'adhésion tardive (MID\_YEAR)}
Si un membre rejoint au mois $m$ (avec $m > 2$) :
$$\texttt{catchUpAmount} = (m - 1) \times (\texttt{sharesCount} \times \texttt{shareUnitAmount} + \texttt{rescueFundContrib}_\text{mensuel})$$
Ce montant doit être versé en intégralité avant que le membre soit éligible à un prêt.
\end{ruleBox}

\subsection{Entité \texttt{ShareCommitment}}

\begin{codeBox}
ShareCommitment {
  id              : ShortUUID(10)  PK
  membershipId    : ShortUUID      FK → Membership  UNIQUE
  sharesCount     : DECIMAL(4,2)   NOT NULL
  monthlyAmount   : DECIMAL(15,2)  NOT NULL  -- figé à l'inscription
  isLocked        : BOOLEAN        DEFAULT false
  lockedAt        : TIMESTAMP      nullable
  lockedById      : ShortUUID FK → User  nullable

  CHECK : sharesCount % 0.25 = 0
  CHECK : sharesCount BETWEEN 0.25 AND 10
}
\end{codeBox}

% -------------------------------------------------------
\section{Module 05 — Sessions et Transactions}

\subsection{Entité \texttt{MonthlySession}}

\begin{codeBox}
MonthlySession {
  id                : ShortUUID(10)   PK
  fiscalYearId      : ShortUUID       FK → FiscalYear
  sessionNumber     : INT             NOT NULL  -- 1..12
  meetingDate       : DATE            NOT NULL
  location          : VARCHAR(500)    nullable
  hostMembershipId  : ShortUUID FK → Membership  nullable
  status            : SessionStatus   DEFAULT DRAFT
  openedAt          : TIMESTAMP       nullable
  openedById        : ShortUUID FK → User  nullable
  closedAt          : TIMESTAMP       nullable
  closedById        : ShortUUID FK → User  nullable
  -- Totaux snapshot --
  totalInscription  : DECIMAL(15,2)   DEFAULT 0
  totalSecours      : DECIMAL(15,2)   DEFAULT 0
  totalCotisation   : DECIMAL(15,2)   DEFAULT 0
  totalPot          : DECIMAL(15,2)   DEFAULT 0
  totalRbtPrincipal : DECIMAL(15,2)   DEFAULT 0
  totalRbtInterest  : DECIMAL(15,2)   DEFAULT 0
  totalEpargne      : DECIMAL(15,2)   DEFAULT 0
  totalProjet       : DECIMAL(15,2)   DEFAULT 0
  totalAutres       : DECIMAL(15,2)   DEFAULT 0
  UNIQUE(fiscalYearId, sessionNumber)
}

SessionStatus : DRAFT → OPEN → REVIEWING → CLOSED
\end{codeBox}

\subsection{Types de Transactions (Entrées)}

\begin{center}
\rowcolors{2}{cayaLight}{white}
\begin{tabular}{L{3.5cm}L{4cm}L{5.5cm}}
  \toprule
  \rowcolor{cayaBlue}
  \textcolor{white}{\textbf{Type}} &
  \textcolor{white}{\textbf{Montant}} &
  \textcolor{white}{\textbf{Description}} \\
  \midrule
  \texttt{INSCRIPTION}    & Configurable         & Frais d'adhésion annuels \\
  \texttt{SECOURS}        & Libre / 50 000 cible & Contribution caisse de secours \\
  \texttt{COTISATION}     & Parts × valeur part  & Cotisation mensuelle tontine \\
  \texttt{POT}            & 3 000 fixe           & Frais de collation réunion \\
  \texttt{RBT\_PRINCIPAL} & Partiel ou total     & Remboursement capital prêt \\
  \texttt{RBT\_INTEREST}  & Calculé              & Remboursement intérêts prêt \\
  \texttt{EPARGNE}        & Libre                & Dépôt d'épargne volontaire \\
  \texttt{PROJET}         & Libre                & Contribution projet spécial \\
  \texttt{AUTRES}         & Libre                & Versements divers \\
  \bottomrule
\end{tabular}
\end{center}

\subsection{Entité \texttt{SessionEntry}}

\begin{codeBox}
SessionEntry {
  id                : ShortUUID(10)     PK
  reference         : VARCHAR(30)       UNIQUE  -- CAYA-2026-03-COT-0042
  sessionId         : ShortUUID FK → MonthlySession  nullable
  membershipId      : ShortUUID         FK → Membership
  type              : TransactionType
  amount            : DECIMAL(15,2)     CHECK (amount > 0)
  loanId            : ShortUUID FK → LoanAccount  nullable
  isOutOfSession    : BOOLEAN           DEFAULT false
  outOfSessionAt    : TIMESTAMP         nullable
  outOfSessionRef   : VARCHAR(100)      nullable
  isImported        : BOOLEAN           DEFAULT false
  notes             : TEXT              nullable
  recordedById      : ShortUUID         FK → User
  recordedAt        : TIMESTAMP         DEFAULT now()
}
\end{codeBox}

\begin{invariantBox}{SESS-01 à SESS-04}
\begin{itemize}
  \item \textbf{SESS-01} : Une \texttt{SessionEntry} normale (\texttt{isOutOfSession = false}) ne peut être créée
        ou modifiée que si \texttt{session.status = OPEN}.
  \item \textbf{SESS-02} : Un remboursement hors-session (\texttt{isOutOfSession = true}) n'accepte que
        les types \texttt{RBT\_PRINCIPAL} ou \texttt{RBT\_INTEREST} et a \texttt{sessionId = null}.
  \item \textbf{SESS-03} : Tout remboursement de prêt génère exactement \textbf{deux entrées} (même \texttt{loanId}) :
        une \texttt{RBT\_INTEREST} et une \texttt{RBT\_PRINCIPAL}.
  \item \textbf{SESS-04} : La session $n$ ne peut être ouverte que si la session $n-1$ est clôturée (ou $n = 1$).
\end{itemize}
\end{invariantBox}

% -------------------------------------------------------
\section{Module 06 — Épargne et Redistribution des Intérêts}

\subsection{Entité \texttt{SavingsLedger}}

\begin{codeBox}
SavingsLedger {
  id                     : ShortUUID(10)  PK
  membershipId           : ShortUUID      FK → Membership  UNIQUE
  balance                : DECIMAL(15,2)  DEFAULT 0
  principalBalance       : DECIMAL(15,2)  DEFAULT 0
  totalInterestReceived  : DECIMAL(15,2)  DEFAULT 0
  lastUpdatedAt          : TIMESTAMP
}
\end{codeBox}

\subsection{Comptes Institutionnels dans le Pool — \texttt{PoolParticipant}}

\begin{codeBox}
PoolParticipant {
  id                     : ShortUUID(10)      PK
  fiscalYearId           : ShortUUID          FK → FiscalYear
  type                   : PoolParticipantType
  label                  : VARCHAR(100)
  initialBalance         : DECIMAL(15,2)      NOT NULL
  currentBalance         : DECIMAL(15,2)      DEFAULT 0
  totalInterestReceived  : DECIMAL(15,2)      DEFAULT 0
  UNIQUE(fiscalYearId, type)
}

PoolParticipantType : RESCUE_FUND | BUREAU
\end{codeBox}

\begin{ruleBox}{Mécanique de Redistribution Mensuelle}
\textbf{Étape 1 — Calcul du pool total d'épargne :}
$$\text{totalSavingsPool} = \sum_i \texttt{SavingsLedger}_i.\texttt{balance} + \sum_j \texttt{PoolParticipant}_j.\texttt{currentBalance}$$

\textbf{Étape 2 — Calcul du pool d'intérêts (selon méthode) :}

Si \texttt{THEORETICAL} :
$$\text{totalInterestPool} = \sum_k \texttt{MonthlyLoanAccrual}_k.\texttt{interestAccrued}$$

Si \texttt{ACTUAL} :
$$\text{totalInterestPool} = \sum \texttt{SessionEntry}[\texttt{RBT\_INTEREST}].\texttt{amount}$$

\textbf{Étape 3 — Allocation proportionnelle :}
$$\texttt{allocationAmount}_i = \frac{\texttt{balance}_i}{\text{totalSavingsPool}} \times \text{totalInterestPool}$$

\textbf{Vérification :} $\sum_i \texttt{allocationAmount}_i = \text{totalInterestPool} \pm 1$ (arrondi)
\end{ruleBox}

% -------------------------------------------------------
\section{Module 07 — Prêts et Remboursements}

\subsection{Entité \texttt{LoanAccount}}

\begin{codeBox}
LoanAccount {
  id                    : ShortUUID(10)  PK
  membershipId          : ShortUUID      FK → Membership
  fiscalYearId          : ShortUUID      FK → FiscalYear
  principalAmount       : DECIMAL(15,2)
  monthlyRate           : DECIMAL(5,4)   DEFAULT 0.04
  disbursedAt           : DATE           nullable
  disbursedById         : ShortUUID FK → User  nullable
  dueBeforeDate         : DATE           NOT NULL
  status                : LoanStatus     DEFAULT PENDING
  outstandingBalance    : DECIMAL(15,2)  DEFAULT 0
  totalInterestAccrued  : DECIMAL(15,2)  DEFAULT 0
  totalRepaid           : DECIMAL(15,2)  DEFAULT 0
  requestedAt           : TIMESTAMP
  requestNotes          : TEXT           nullable
}

LoanStatus : PENDING | ACTIVE | PARTIALLY_REPAID | CLOSED
\end{codeBox}

\subsection{Entité \texttt{MonthlyLoanAccrual} — Calcul de l'Encours Glissant}

\begin{codeBox}
MonthlyLoanAccrual {
  id                    : ShortUUID(10)  PK
  loanId                : ShortUUID      FK → LoanAccount
  sessionId             : ShortUUID      FK → MonthlySession
  month                 : INT
  balanceAtMonthStart   : DECIMAL(15,2)  -- encours[m-1]
  interestAccrued       : DECIMAL(15,2)  -- = balanceAtMonthStart × 0.04
  balanceWithInterest   : DECIMAL(15,2)  -- = balanceAtMonthStart × 1.04
  repaymentReceived     : DECIMAL(15,2)  DEFAULT 0
  balanceAtMonthEnd     : DECIMAL(15,2)  -- = balanceWithInterest - repayment
  UNIQUE(loanId, sessionId)
  CHECK : balanceAtMonthEnd >= 0
}
\end{codeBox}

\begin{ruleBox}{Algorithme de Calcul des Intérêts (Encours Glissant)}
Soit $P_0$ le capital initial emprunté au mois 1 :

\begin{align*}
  \text{Mois 1} &: \text{encours}_1 = P_0 \\
  \text{Mois 2} &: \text{dû}_2 = \text{encours}_1 \times 1{,}04 \\
                &\quad \text{encours}_2 = \text{dû}_2 - R_2 \\
  \text{Mois } m &: \text{dû}_m = \text{encours}_{m-1} \times 1{,}04 \\
                 &\quad \text{encours}_m = \text{dû}_m - R_m
\end{align*}

\textbf{Règle de décomposition d'un remboursement $R$ :}
\begin{itemize}
  \item Si $R \geq \text{interestDue}$ : intérêtPart = interestDue ; capitalPart = $R -$ interestDue
  \item Si $R < \text{interestDue}$ : tout va aux intérêts ; capital inchangé ; solde = $\text{dû} - R$
\end{itemize}
\end{ruleBox}

\subsection{Conditions d'Éligibilité à un Prêt}

\begin{center}
\rowcolors{2}{cayaLight}{white}
\begin{tabular}{L{6cm}R{6cm}}
  \toprule
  \rowcolor{cayaBlue}
  \textcolor{white}{\textbf{Condition}} &
  \textcolor{white}{\textbf{Valeur}} \\
  \midrule
  Épargne minimale requise         & 100 000 XAF \\
  Ratio prêt / épargne maximum     & 5 (ex: 100k épargne → 500k max) \\
  Montant minimum                  & 10 000 XAF \\
  Montant maximum absolu           & 1 000 000 XAF \\
  Prêts simultanés autorisés       & 2 maximum \\
  Date d'échéance maximale         & avant cassation \\
  Pénalité de retard               & aucune \\
  \bottomrule
\end{tabular}
\end{center}

\begin{invariantBox}{LOAN-01 à LOAN-05}
\begin{itemize}
  \item \textbf{LOAN-01} : Imputation FIFO en cas de prêts multiples — les remboursements sont affectés
        au prêt le plus ancien (\texttt{requestedAt ASC}).
  \item \textbf{LOAN-02} : Un prêt non soldé à la cassation génère un \texttt{CarryoverLoanRecord}
        avec \texttt{carryoverAmount = outstandingBalance × 1.04} et une notification dans l'application.
  \item \textbf{LOAN-03} : Le statut passe \texttt{CLOSED} uniquement si \texttt{outstandingBalance = 0}.
  \item \textbf{LOAN-04} : Un remboursement inférieur aux intérêts dus est accepté ; les intérêts
        sont capitalisés dans le solde restant.
  \item \textbf{LOAN-05} : \texttt{dueBeforeDate < FiscalYear.cassationDate}.
\end{itemize}
\end{invariantBox}

% -------------------------------------------------------
\section{Module 08 — Caisse de Secours}

\subsection{Entités}

\begin{codeBox}
RescueFundLedger {
  id                : ShortUUID(10)  PK
  fiscalYearId      : ShortUUID      FK → FiscalYear  UNIQUE
  totalBalance      : DECIMAL(15,2)  DEFAULT 0
  targetPerMember   : DECIMAL(15,2)  DEFAULT 50000
  minimumPerMember  : DECIMAL(15,2)  DEFAULT 25000
  memberCount       : INT
  targetTotal       : DECIMAL(15,2)  -- calculé = target × memberCount
}

RescueFundEvent {
  id              : ShortUUID(10)  PK
  ledgerId        : ShortUUID      FK → RescueFundLedger
  beneficiaryId   : ShortUUID      FK → Membership
  eventType       : RescueEventType
  amount          : DECIMAL(15,2)
  authorizedById  : ShortUUID      FK → User
  eventDate       : DATE
  description     : TEXT           nullable
  createdAt       : TIMESTAMP
}

RescueFundPosition {
  id              : ShortUUID(10)  PK
  membershipId    : ShortUUID      FK → Membership  UNIQUE
  ledgerId        : ShortUUID      FK → RescueFundLedger
  paidAmount      : DECIMAL(15,2)  DEFAULT 0
  balance         : DECIMAL(15,2)  DEFAULT 0
  refillDebt      : DECIMAL(15,2)  DEFAULT 0
}
\end{codeBox}

\begin{invariantBox}{RES-01 à RES-03}
\begin{itemize}
  \item \textbf{RES-01} : Après tout versement de secours, \texttt{totalBalance $\geq$ minimumPerMember × memberCount}.
  \item \textbf{RES-02} : Le montant d'un événement est fixé par \texttt{RescueEventAmount[eventType]} — non négociable.
  \item \textbf{RES-03} : L'autorisation d'un versement requiert \texttt{role ∈ \{PRESIDENT, VICE\_PRESIDENT\}}.
\end{itemize}
\end{invariantBox}

% -------------------------------------------------------
\section{Module 09 — Bénéficiaires}

\begin{codeBox}
BeneficiarySchedule {
  id            : ShortUUID(10)  PK
  fiscalYearId  : ShortUUID      FK → FiscalYear  UNIQUE
  createdAt     : TIMESTAMP
}

BeneficiarySlot {
  id               : ShortUUID(10)       PK
  scheduleId       : ShortUUID           FK → BeneficiarySchedule
  sessionId        : ShortUUID           FK → MonthlySession
  month            : INT
  slotIndex        : INT                 -- permet plusieurs bénéficiaires/session
  membershipId     : ShortUUID FK → Membership  nullable
  amountDelivered  : DECIMAL(15,2)       DEFAULT 0
  designatedById   : ShortUUID FK → User nullable
  designatedAt     : TIMESTAMP           nullable
  status           : BeneficiaryStatus   DEFAULT UNASSIGNED
  deliveredAt      : TIMESTAMP           nullable
  notes            : TEXT                nullable
  UNIQUE(sessionId, slotIndex)
  UNIQUE(sessionId, membershipId)
}

BeneficiaryStatus : UNASSIGNED | ASSIGNED | DELIVERED
\end{codeBox}

\begin{ruleBox}{Règles de désignation des bénéficiaires}
\begin{itemize}
  \item La désignation est \textbf{manuelle} par \texttt{PRESIDENT} ou \texttt{VICE\_PRESIDENT}.
  \item Le montant remis est \textbf{saisi manuellement} par le Président/VP.
  \item Plusieurs slots par session sont autorisés (\texttt{slotIndex} 0, 1, 2\ldots).
  \item Un membre ne peut pas recevoir deux slots dans la même session.
  \item Le statut \texttt{DELIVERED} n'est possible que si la session est \texttt{CLOSED}.
\end{itemize}
\end{ruleBox}

% -------------------------------------------------------
\section{Module 10 — Cassation et Rollover}

\begin{codeBox}
CassationRecord {
  id                    : ShortUUID(10)  PK
  fiscalYearId          : ShortUUID      FK → FiscalYear  UNIQUE
  executedAt            : TIMESTAMP
  executedById          : ShortUUID      FK → User
  totalSavingsReturned  : DECIMAL(15,2)
  totalInterestReturned : DECIMAL(15,2)
  totalDistributed      : DECIMAL(15,2)
  memberCount           : INT
  notes                 : TEXT
}

CassationRedistribution {
  id              : ShortUUID(10)  PK
  cassationId     : ShortUUID      FK → CassationRecord
  membershipId    : ShortUUID      FK → Membership
  savingsAmount   : DECIMAL(15,2)
  interestAmount  : DECIMAL(15,2)
  totalReturned   : DECIMAL(15,2)
  UNIQUE(cassationId, membershipId)
}

CarryoverLoanRecord {
  id               : ShortUUID(10)  PK
  originalLoanId   : ShortUUID      FK → LoanAccount  UNIQUE
  newFiscalYearId  : ShortUUID      FK → FiscalYear
  carryoverAmount  : DECIMAL(15,2)  -- = outstandingBalance × 1.04
  reason           : TEXT
  approvedById     : ShortUUID      FK → User
  createdAt        : TIMESTAMP
}
\end{codeBox}

\begin{ruleBox}{Processus de Cassation (séquence ordonnée)}
\begin{enumerate}
  \item Identifier tous les \texttt{LoanAccount} avec \texttt{status ≠ CLOSED}
        → créer un \texttt{CarryoverLoanRecord} par prêt actif + notification dans l'application.
  \item Exécuter la distribution d'intérêts de la dernière session.
  \item Calculer \texttt{CassationRedistribution} pour chaque membre actif.
  \item Distribuer les soldes \texttt{PoolParticipant} (RESCUE\_FUND + BUREAU).
  \item Passer \texttt{FiscalYear.status = CLOSED}.
  \item Créer \texttt{CassationRecord} avec tous les totaux.
  \item Notifier PRESIDENT et TRESORIER via WhatsApp/SMS.
  \item Initialiser le \texttt{FiscalYear} N+1 en statut \texttt{PENDING}.
\end{enumerate}
\end{ruleBox}

% ============================================================
\chapter{Fonctionnalités par Module}
% ============================================================

\section{Gestion des Sessions Mensuelles}

\begin{ruleBox}{Cycle de vie d'une session}
\textbf{DRAFT} : La session est créée automatiquement lors de l'activation de l'exercice (12 sessions).\\
\textbf{OPEN} : Le Trésorier ouvre la session le jour de la réunion. Les entrées sont saisies.\\
\textbf{REVIEWING} : Le Trésorier clôt la saisie pour revue par le Président.\\
\textbf{CLOSED} : Le Président valide. Données figées. Intérêts distribués automatiquement.
\end{ruleBox}

\subsection{Saisie des Transactions}

Lors d'une session \texttt{OPEN}, le Trésorier saisit membre par membre :
\begin{itemize}
  \item Pour \textbf{chaque entrée}, une référence unique est générée au format :
        \texttt{CAYA-AAAA-MM-TYPE-NNNN}
  \item Pour un \textbf{remboursement de prêt}, le système décompose automatiquement le montant
        en intérêts dus et capital selon la règle de l'encours glissant.
  \item Les \textbf{remboursements hors-session} (Mobile Money / Orange Money) sont saisis
        directement sans session ouverte, avec une référence de virement.
\end{itemize}

\section{Rapports et Exports}

\subsection{Tableau Récapitulatif Annuel (équivalent Rem25)}

Ce rapport correspond exactement à la feuille \textit{Rem25} du fichier Excel, mais entièrement calculé :

\begin{center}
\small
\rowcolors{2}{cayaLight}{white}
\begin{tabular}{L{4cm}L{9cm}}
  \toprule
  \rowcolor{cayaBlue}
  \textcolor{white}{\textbf{Colonne}} &
  \textcolor{white}{\textbf{Source de calcul}} \\
  \midrule
  N° / Nom          & \texttt{MemberProfile.memberCode + fullName()} \\
  Inscription       & $\sum$ \texttt{SessionEntry[INSCRIPTION]} \\
  Secours           & $\sum$ \texttt{SessionEntry[SECOURS]} \\
  Cotisation        & $\sum$ \texttt{SessionEntry[COTISATION]} \\
  Pot               & $\sum$ \texttt{SessionEntry[POT]} \\
  Épargne brute     & $\sum$ \texttt{SavingsEntry[DEPOSIT].amount} \\
  Intérêts reçus    & \texttt{SavingsLedger.totalInterestReceived} \\
  NAP (Épargne+Int) & \texttt{SavingsLedger.balance} \\
  Prêts totaux      & $\sum$ \texttt{LoanAccount.principalAmount} \\
  Intérêts dus      & $\sum$ \texttt{MonthlyLoanAccrual.interestAccrued} \\
  Remboursé         & $\sum$ \texttt{LoanRepayment.amount} \\
  Reste à payer     & \texttt{outstandingBalance × 1.04} \\
  \bottomrule
\end{tabular}
\end{center}

\subsection{Formats d'Export}

\begin{itemize}
  \item \textbf{PDF} : En-tête CAYA, logo, tous rapports et fiches individuelles.
  \item \textbf{Excel (.xlsx)} : Tableau récapitulatif annuel, données mensuelles.
  \item \textbf{Fiche individuelle} : Vue détaillée d'un membre (PDF + Excel).
\end{itemize}

% ============================================================
\chapter{Exigences Non Fonctionnelles}
% ============================================================

\section{Performance}

\begin{itemize}
  \item La table \texttt{AuditLog} est \textbf{partitionnée par année} via \texttt{RANGE PARTITIONING}
        MySQL 8.0 pour maintenir des performances constantes après 5 ans d'utilisation.
  \item Index obligatoires :
        \texttt{SessionEntry(membershipId, type)},
        \texttt{MonthlyLoanAccrual(sessionId)},
        \texttt{LoanAccount(membershipId, status)}.
\end{itemize}

\section{Sécurité}

\begin{itemize}
  \item Authentification par \textbf{phone ou username} (pas d'email). JWT + Refresh Token.
  \item Toute action sensible (modification config, forçage SUPER\_ADMIN, réouverture session)
        génère une entrée \textbf{AuditLog immuable} (\textit{append-only}).
  \item L'audit log n'est lisible que par \texttt{SUPER\_ADMIN}, \texttt{PRESIDENT} et \texttt{SECRETAIRE\_GENERAL}.
  \item Les photos membres sont stockées en \textbf{BYTEA} (binaire, max 100 KB) — pas d'URLs externes.
\end{itemize}

\section{Intégrité des Données}

\begin{itemize}
  \item Toutes les valeurs financières utilisent le type \texttt{DECIMAL(15,2)} — zéro flottant
        (\textbf{pas} \texttt{FLOAT} ni \texttt{DOUBLE} sous MySQL).
  \item Le trigger anti-chevauchement d'exercices est implémenté au niveau base de données.
  \item La table \texttt{TransactionSequence} garantit l'unicité des références de transactions
        par un verrou atomique (\texttt{SELECT FOR UPDATE}).
  \item Aucune valeur financière calculable n'est saisie manuellement (principe de vue calculée).
\end{itemize}

\section{Adaptations spécifiques MySQL 8.0}

MySQL présente plusieurs différences par rapport à PostgreSQL qui impactent l'implémentation :

\begin{center}
\small
\rowcolors{2}{cayaLight}{white}
\begin{tabular}{L{3.5cm}L{4.5cm}L{5cm}}
  \toprule
  \rowcolor{cayaBlue}
  \textcolor{white}{\textbf{Besoin}} &
  \textcolor{white}{\textbf{Solution MySQL 8.0}} &
  \textcolor{white}{\textbf{Note}} \\
  \midrule
  Types monétaires        & \texttt{DECIMAL(15,2)}     & Jamais FLOAT ni DOUBLE \\
  ENUM métier             & \texttt{ENUM(...)}          & Natif MySQL \\
  JSON (audit before/after)& \texttt{JSON}              & Supporté depuis MySQL 5.7 \\
  Partitionnement AuditLog & \texttt{RANGE} sur \texttt{YEAR(createdAt)} & Requiert moteur InnoDB \\
  Trigger BEFORE INSERT   & Trigger MySQL natif         & Anti-chevauchement exercices \\
  Verrou séquence         & \texttt{SELECT ... FOR UPDATE} & Natif InnoDB \\
  UUID courts             & \texttt{VARCHAR(10)} + \texttt{UUID()} tronqué ou Nanoid & Généré applicatif \\
  Vues calculées          & \texttt{CREATE VIEW}        & Rapport Rem25, LoanSummary \\
  Rollback transactionnel & InnoDB uniquement           & Moteur obligatoire (pas MyISAM) \\
  \bottomrule
\end{tabular}
\end{center}

\begin{warningBox}
  \textbf{Moteur InnoDB obligatoire} sur toutes les tables. MyISAM ne supporte pas les transactions
  ni les clés étrangères — il ne doit \textbf{jamais} être utilisé dans ce projet.
  Ajouter \texttt{ENGINE=InnoDB DEFAULT CHARSET=utf8mb4 COLLATE=utf8mb4\_unicode\_ci}
  sur chaque \texttt{CREATE TABLE}.
\end{warningBox}

\section{Architecture Flutter (Mobile)}

L'application mobile Flutter couvre les fonctionnalités membres :

\begin{itemize}
  \item Consultation de la fiche personnelle (épargne, prêts, secours, cotisations).
  \item Historique des transactions de l'exercice en cours.
  \item Notifications push (rappels réunion, confirmation paiement, alerte remboursement).
  \item Visualisation des graphiques d'évolution (épargne + intérêts cumulés).
  \item Export PDF de la fiche individuelle.
\end{itemize}

\begin{center}
\small
\rowcolors{2}{cayaLight}{white}
\begin{tabular}{L{3.5cm}L{4cm}L{5.5cm}}
  \toprule
  \rowcolor{cayaBlue}
  \textcolor{white}{\textbf{Couche Flutter}} &
  \textcolor{white}{\textbf{Package recommandé}} &
  \textcolor{white}{\textbf{Usage}} \\
  \midrule
  State management  & Riverpod / Bloc     & Gestion état applicatif \\
  HTTP client       & Dio                 & Appels API NestJS \\
  Auth tokens       & flutter\_secure\_storage & Stockage JWT sécurisé \\
  Graphiques        & fl\_chart           & Courbes épargne/intérêts \\
  PDF export        & pdf + printing      & Fiche membre \\
  Notifications push& firebase\_messaging & Rappels, confirmations \\
  Navigation        & go\_router          & Routing déclaratif \\
  \bottomrule
\end{tabular}
\end{center}

\section{Portabilité et Accessibilité}

\begin{itemize}
  \item \textbf{Interface web} (Next.js) : Trésorier, Bureau — saisie complète.
  \item \textbf{Application mobile} (Flutter / Dart) : Membres — consultation, notifications.
        Une seule codebase compilée nativement pour Android et iOS.
  \item \textbf{Notifications} : WhatsApp Business API + SMS (3 jours avant réunion,
        confirmation de paiement, rappel de remboursement).
\end{itemize}

\section{Évolutivité}

L'architecture modulaire garantit que chaque module peut être développé, testé, et déployé
indépendamment. La conception intègre les points de variabilité suivants via configuration :
\begin{itemize}
  \item Méthode de calcul des intérêts (\texttt{THEORETICAL / ACTUAL}) — switchable par exercice.
  \item Dates d'exercice — entièrement configurables par le Président.
  \item Montants de secours — reconfigurables chaque exercice.
  \item Frais d'inscription — champ distinct nouveau / retour.
\end{itemize}

% ============================================================
\chapter{Plan de Développement Modulaire}
% ============================================================

\section{Ordre de Développement Recommandé}

\begin{center}
\rowcolors{2}{cayaLight}{white}
\begin{tabular}{C{0.8cm}L{3.5cm}L{4cm}L{4.5cm}}
  \toprule
  \rowcolor{cayaBlue}
  \textcolor{white}{\textbf{Sprint}} &
  \textcolor{white}{\textbf{Module}} &
  \textcolor{white}{\textbf{Livrable}} &
  \textcolor{white}{\textbf{Dépendances}} \\
  \midrule
  1 & Auth \& RBAC            & Login, rôles, JWT            & Aucune \\
  2 & Configuration           & Paramètres, snapshot         & Auth \\
  3 & Membres                 & Profils, codes, photos       & Auth \\
  4 & Exercice Fiscal         & FiscalYear, Membership       & Config, Membres \\
  5 & Sessions \& Transactions& Saisie, références, audit    & Exercice \\
  6 & Épargne \& Intérêts     & Ledgers, redistribution      & Sessions \\
  7 & Prêts                   & Octroi, accrual, rembt       & Sessions, Épargne \\
  8 & Caisse de Secours       & Fonds, événements            & Membres, Exercice \\
  9 & Bénéficiaires           & Planning, désignation        & Sessions \\
 10 & Cassation \& Rollover   & Clôture, report prêts        & Tous \\
 11 & Rapports \& Exports     & PDF, Excel, graphiques       & Tous \\
 12 & Notifications           & WhatsApp, SMS                & Tous \\
  \bottomrule
\end{tabular}
\end{center}

\section{Critères de Validation par Module}

Chaque module est considéré \textbf{validé} lorsque :
\begin{enumerate}
  \item Tous les invariants métier de la section sont couverts par des tests automatisés.
  \item Aucune valeur financière calculable n'est saisie en dur.
  \item Le module peut être déployé sans impacter les modules déjà en production.
  \item Un scénario de bout en bout (happy path + cas limites) est documenté et passant.
\end{enumerate}

% ============================================================
\chapter*{Annexes}
\addcontentsline{toc}{chapter}{Annexes}
% ============================================================

\section*{A. Glossaire}

\begin{description}
  \item[Part] Unité de cotisation mensuelle d'un membre (100 000 XAF). Fractions autorisées : demi-part,
              quart de part. Engagement figé pour toute la durée de l'exercice.
  \item[Cassation] Clôture de l'exercice : redistribution de l'épargne accumulée et des intérêts générés
                   à chaque membre proportionnellement à son épargne.
  \item[Encours] Solde restant dû d'un prêt à un instant donné, après application des intérêts et déduction
                 des remboursements.
  \item[Taux glissant] Méthode de calcul des intérêts : chaque mois, l'intérêt est calculé sur l'encours
                       réel du mois précédent, pas sur le capital initial.
  \item[Pool d'intérêts] Montant total des intérêts générés par les prêts actifs durant un mois, redistribué
                          proportionnellement aux épargnants (membres + comptes institutionnels).
  \item[NAP] \textit{Net À Payer} — montant que la caisse doit restituer à un membre à la cassation :
             épargne totale $+$ intérêts accumulés.
  \item[Pot] Cotisation mensuelle de 3 000 XAF par membre pour financer la collation lors de la réunion.
             Comptabilisé séparément, n'entre pas dans le pool d'épargne.
  \item[Rattrapage (catch-up)] Montant dû par un membre arrivé en cours d'exercice pour couvrir les
                               cotisations des mois précédents.
\end{description}

\section*{B. Récapitulatif des Invariants Métier}

\begin{center}
\small
\rowcolors{2}{cayaLight}{white}
\begin{tabular}{L{2cm}L{11cm}}
  \toprule
  \rowcolor{cayaBlue}
  \textcolor{white}{\textbf{Code}} &
  \textcolor{white}{\textbf{Règle}} \\
  \midrule
  AUTH-01  & Un seul PRESIDENT actif simultanément \\
  AUTH-02  & Un seul VICE\_PRESIDENT actif simultanément \\
  AUTH-03  & phone User synchronisé avec MemberProfile.phone1 \\
  AUTH-04  & Tout changement de rôle tracé dans AuditLog \\
  CONF-01  & FiscalYearConfig immuable après activation de l'exercice \\
  CONF-02  & Modification en cours d'exercice requiert SUPER\_ADMIN + motif \\
  CONF-03  & interestPoolMethod switchable uniquement en PENDING \\
  SESS-01  & SessionEntry créée uniquement si session.status = OPEN \\
  SESS-02  & Remboursement hors-session = RBT\_PRINCIPAL ou RBT\_INTEREST uniquement \\
  SESS-03  & Remboursement prêt = 2 entries obligatoires (interest + principal) \\
  SESS-04  & Session N ne s'ouvre que si session N-1 est CLOSED \\
  LOAN-01  & Imputation remboursements FIFO par ancienneté du prêt \\
  LOAN-02  & Prêt non soldé à cassation → CarryoverLoanRecord + notification \\
  LOAN-03  & LoanAccount.status = CLOSED iff outstandingBalance = 0 \\
  LOAN-04  & Remboursement < intérêts dus → capitalisé dans le solde \\
  LOAN-05  & dueBeforeDate < FiscalYear.cassationDate \\
  RES-01   & Solde caisse de secours ≥ minimumPerMember × memberCount après versement \\
  RES-02   & Montant secours = valeur fixe par type d'événement \\
  RES-03   & Autorisation secours requiert PRESIDENT ou VICE\_PRESIDENT \\
  \bottomrule
\end{tabular}
\end{center}

\section*{C. Table de Correspondance Excel → Système}

\begin{center}
\small
\rowcolors{2}{cayaLight}{white}
\begin{tabular}{L{3.5cm}L{4cm}L{5.5cm}}
  \toprule
  \rowcolor{cayaBlue}
  \textcolor{white}{\textbf{Feuille Excel}} &
  \textcolor{white}{\textbf{Entité/Vue}} &
  \textcolor{white}{\textbf{Remarque}} \\
  \midrule
  \texttt{ep+int (ancien)}  & \texttt{InterestDistributionSnapshot} & Méthode THEORETICAL \\
  \texttt{ep+int (2)}       & \texttt{InterestDistributionSnapshot} & Méthode ACTUAL \\
  \texttt{prets}            & \texttt{LoanAccount} (encours cumulé) & Vue calculée \\
  \texttt{prets25}          & \texttt{LoanAccount} (nouveaux prêts) & Vue calculée \\
  \texttt{intérêts25}       & \texttt{MonthlyLoanAccrual}           & Automatisé \\
  \texttt{intérêts}         & \texttt{MonthlyLoanAccrual}           & Automatisé \\
  \texttt{Rem25}            & Vue \texttt{LoanSummaryView}          & Zéro saisie manuelle \\
  \texttt{insc+sec}         & \texttt{SessionEntry} (INSCRIPTION, SECOURS) & Tracé par session \\
  \bottomrule
\end{tabular}
\end{center}

\vfill
\begin{center}
  \textcolor{cayaBlue}{\rule{0.5\textwidth}{0.4pt}}\\[0.3cm]
  {\small\textit{Document préparé par l'équipe d'Ingénierie Logicielle CAYA}}\\
  {\small\textcolor{cayaGray!50!black}{Version 1.0 — Mars 2026 — Confidentiel}}
\end{center}

\end{document}
